% =========================
% Chapter 2
% =========================
\chapter{Background}
\label{ch:background}

This chapter introduces the technical building blocks Lumen is assembled from: visual
perception (object detection, spatial reasoning, segmentation, and depth), the voice
pipeline (speech recognition, command understanding, speech synthesis), device
sensing available to a web browser, and state-based system control.

\section{Vision-Based Perception}

\subsection{Object Detection}

Vision-based object detection enables Lumen to perceive and interpret the surrounding
environment by identifying relevant objects and estimating their spatial location
from camera input. Modern object detectors process an input image to simultaneously
predict object class labels and corresponding bounding boxes, allowing real-time
localization and recognition within a single inference
pipeline~\cite{girshick2015fast}. Among these approaches, one-stage detectors such as
the You Only Look Once (YOLO) family are particularly suitable for assistive systems
due to their low inference latency and unified detection formulation, which directly
maps images to object predictions without intermediate proposal
stages~\cite{redmon2016yolo}. The YOLOv8 family~\cite{jocher2023yolov8} spans a range
of model sizes trading accuracy for speed, all pretrained on the COCO
dataset~\cite{lin2014coco}, whose 80 everyday object classes (cup, chair,
refrigerator, bed, \ldots) align well with the vocabulary of indoor assistive tasks.

Lumen uses two YOLOv8 variants for two different jobs. \textbf{YOLOv8n} (nano, ${\sim}6$\,MB)
serves Object Allocation, where the target classes are large in frame and the
detection loop must run at ${\sim}5$\,FPS even on CPU. \textbf{YOLOv8m} (medium)
serves Navigation, where the decisive factor is recall on room-indicator objects seen
at an angle and across a room --- detection accuracy, not inference time, is the
bottleneck there. In addition, Navigation uses two custom-trained YOLOv8 door models
(Section~\ref{sec:door_model}), since COCO contains no door class.

\subsection{Relative Spatial Reasoning}

Bounding boxes alone do not tell a blind user anything actionable; they must be
translated into egocentric guidance. Lumen infers the target's horizontal
\emph{region} (left / center / right) from the box centre, and a coarse
\emph{distance} bucket (near / medium / far) from the box's apparent size. Because
apparent size conflates true size with distance, Lumen applies \emph{per-class size
priors}: a laptop occupying 40\% of the frame width is near, and a cup occupying 18\%
is equally near, because cups are small. Qualitative spatial representations of this
kind have been shown sufficient for indoor assistive guidance, where approximate
awareness --- not metric depth --- drives safe and intuitive
instructions~\cite{liu2019spatial}. Where metric distance genuinely matters --- how
many steps to a door --- Lumen instead uses a pinhole-camera model over a known
object height (Section~\ref{sec:door_distance}).

\subsection{Semantic Segmentation and Monocular Depth}

Object detectors can only report objects they have a name for, yet the clutter that
actually litters floors --- clothes piles, boxes, bins --- is largely absent from
detection vocabularies. Two class-agnostic perception tools close this gap.
\emph{Semantic segmentation} labels every pixel with a scene category; transformer
architectures such as SegFormer~\cite{xie2021segformer}, trained on the ADE20K scene
dataset~\cite{zhou2017ade20k}, can classify floor, rug, wall, and door pixels in real
time. This turns the question ``is there an obstacle?'' into the more tractable ``is
the strip of floor I am about to walk on still floor?''. \emph{Monocular depth
estimation} models such as Depth Anything V2~\cite{yang2024depthanything} predict
relative per-pixel depth from a single image, enabling a complementary tripwire
signal: a patch of the walking lane that reads much nearer than the floor beside it
is an intrusion. Lumen's obstacle watchdog (Section~\ref{sec:obstacles}) uses floor
segmentation as its primary class-agnostic signal, with the depth tripwire retained
as an alternative.

\subsection{Hand Pose Estimation}

Reach Guidance requires knowing where the user's hand --- specifically the fingertip
--- is relative to the target object. MediaPipe Hands~\cite{zhang2020mediapipe}
provides 21 three-dimensional hand landmarks from a monocular camera in real time on
CPU, with models compact enough to ship inside the library itself. Lumen uses the
index fingertip and wrist landmarks to compute hand-to-target guidance and to detect
the moment the fingertip enters the target's bounding box.

\section{Speech-to-Text (STT)}

Speech-to-Text serves as the primary input modality in voice-driven assistive
systems, converting spoken commands into text for downstream interpretation. This
enables hands-free interaction, which is essential for users who cannot rely on
visual interfaces~\cite{graves2013speech}. A typical STT pipeline consists of audio
capture, acoustic feature extraction, and neural
transcription~\cite{hinton2012deep}.

Lumen adopts Whisper~\cite{radford2022whisper} as its recognition model, in the
CTranslate2-based \emph{faster-whisper} implementation, which runs the same weights
several times faster on CPU. Whisper is trained on large-scale, diverse audio,
yielding robust transcription across accents, speaking styles, and background noise
--- the conditions of real households. Robustness matters more than vocabulary here:
Lumen's utterances are short imperative commands, and the transcription errors that
do occur (``find my cop'' for \emph{cup}, ``fone'' for \emph{phone}) are handled
downstream by the command parser rather than by widening the recognition
model~\cite{li2017robust}.

\section{Language Understanding}
\label{sec:command_parsing}

After transcription, the text must be mapped to an explicit system action. Large
language models offer great flexibility for this mapping~\cite{brown2020language},
but assistive systems demand deterministic, low-latency, and auditable behavior;
misinterpreting a command has physical consequences. Lumen therefore uses a
constrained command parser built on fuzzy string matching: every possible
(prefix, target) split of the utterance is scored against known command templates
(``find my \ldots'', ``navigate to the \ldots'') and known target vocabularies, with
synonym and mishearing normalization applied first. This follows the long-standing
insight from task-oriented dialog systems that a limited, well-defined intent
schema yields reliable and predictable interaction~\cite{tur2011spoken,
williams2013dialog}. Completion phrases (``got it'', ``stop'') and information
queries (``describe'', ``repeat that'') are recognized the same way during active
tasks.

\section{Text-to-Speech (TTS)}

Text-to-Speech completes the closed interaction loop, delivering guidance in auditory
form --- the primary output channel for blind users~\cite{preece2015interaction}.
Modern neural pipelines synthesize natural speech from text in a single
architecture~\cite{wang2017tacotron}; classical parametric synthesis research
established the clarity and latency requirements that interactive systems must
meet~\cite{zen2009parametric}. Lumen uses Google's gTTS service, which produces
consistent, natural MP3 audio~\cite{googlecloudtts}. Two engineering measures adapt a
cloud TTS service to real-time guidance: synthesized clips are cached in a bounded
LRU keyed by exact text (guidance phrases repeat constantly), and synthesis runs in a
background speech queue so that network latency never stalls perception
(Section~\ref{sec:speech_pipeline}). Assistive TTS additionally demands
\emph{discipline} over \emph{quantity}: Lumen enforces a two-tier speech policy in
which priority announcements always play, while ambient nudges are dropped rather
than queued whenever the audio pipe is busy.

\section{Browser Device Sensing}

Delivering the client as a web page constrains sensing to what browsers expose ---
and modern browsers expose enough. \texttt{getUserMedia} provides camera frames and
microphone audio (requiring HTTPS outside localhost); \texttt{MediaRecorder}
captures push-to-talk utterances; \texttt{DeviceMotionEvent} reports coarse
accelerometer activity, which Lumen classifies into still / moving / walking;
\texttt{DeviceOrientationEvent} reports the compass heading that anchors the
navigation room scan; the Vibration API delivers haptic cues; and the Wake Lock API
keeps the screen (and thus the camera) alive. iOS Safari imposes additional
constraints --- sensor permissions must be requested from within a user gesture, and
audio playback must be unlocked by a user action --- which shaped the client's
startup sequence (Section~\ref{sec:frontend}).

\section{State-Based System Control}

Finite state machines provide deterministic behavior, clear task boundaries, and
predictable recovery --- properties long exploited in robotics control
architectures~\cite{brooks1986robust} and their modern successors such as behavior
trees~\cite{colledanchise2018behavior}. In Lumen, an explicit FSM is the authority on
what the system is doing at every moment: idle, listening for a command, executing a
specific task, or returning to idle. Perception runs only inside active task states,
which is precisely what implements the task-scoped paradigm; the FSM's explicit
transitions are also what make the system testable, since every behavioral claim can
be phrased as ``event X in state Y produces state Z and utterance W''. The
exploration navigation task nests a second, task-internal state machine (scan, face
target, approach door, confirm arrival) beneath the session FSM, following the same
principle at a finer grain.
